\section*{Chapter \ref{ch:ml:representation-learning} --- Representation Learning}

\begin{solution}{exo:ml:representation-learning:1}
Among $2m - 1 = 511$ candidates; with uninformative codes every candidate is equally likely and the loss is
$\ln 511 = 6.24$.
\end{solution}

\begin{solution}{exo:ml:representation-learning:2}
$300\times0.5 = 150$ seconds, 2.5 minutes; $4\,444\times0.5 = 2\,222$ seconds, 37.0 minutes.
\end{solution}

\begin{solution}{exo:ml:representation-learning:3}
$400\times64 + 64 + 64\times16 + 16 = 26\,704$ parameters; a \omterm{def:ml:representation-learning:probe}{linear probe} on the code has $16 + 1 = 17$.
\end{solution}

\begin{solution}{exo:ml:representation-learning:4}
A reconstruction objective keeps the directions of largest variance (\cref{prop:ml:represent:pca} for the linear case),
and on an order book those are the deep levels' sizes; the top-of-book imbalance that predicts the next seconds is a small
share of the variance and is partly lost in a 16-dimensional code. The raw window keeps it, and ridge can find it.
\end{solution}

\begin{solution}{exo:ml:representation-learning:5}
Ridge on 400 raw inputs has many parameters to estimate; with few labels its estimates are noisy and a 16-dimensional
code that already points at the right directions wins. With 4\,444 labels ridge estimates the useful directions itself,
and the code's compression starts to cost more than it saves.
\end{solution}

\begin{solution}{exo:ml:representation-learning:6}
\omterm{def:ml:representation-learning:ssl}{Pre-training} on the test months lets the encoder fit their particular states: a transductive leak that no label is needed
for (distribution shift in the test period is absorbed into the representation). Pre-train only on data before the test
period, as for any fitted component.
\end{solution}

\begin{solution}{exo:ml:representation-learning:7}
\texttt{compare(1000, ahead=16)}: the predictive probe falls to 0.204, from 0.285 at two seconds, level with the network
from scratch. The book eight seconds ahead is mostly unpredictable, so the pretext's target is mostly noise and the
encoder learns less from it.
\end{solution}

\begin{solution}{exo:ml:representation-learning:8}
For view $a$ with partner $p(a)$, the loss term is $-\ln\bigl(e^{s_{a,p(a)}}/\sum_{b\ne a}e^{s_{ab}}\bigr)$: the
cross-entropy of a softmax over the $2m - 1$ other views with the partner as the correct class. If all similarities are
equal each term is $-\ln(1/(2m-1)) = \ln(2m - 1)$, and so is their average.
\end{solution}

\begin{solution}{pb:ml:representation-learning:1}
\begin{enumerate}
\item Reconstruct a masked, noisy window; match two augmented views of a window against other windows; forecast the change
of the book two seconds after the window.
\item The forecasting pretext: its target is the stream's own near future, the same kind of object as the task's.
\item That its code keeps the directions of largest variance, which need not be predictive.
\item That Gaussian noise and random masking of entries do not change what matters about a window.
\item 0.22, 0.12 and 0.18.
\item 0.29, 0.24 and 0.20.
\item Worse than the raw window at 1\,000 and 4\,444; at 300 the \omterm{def:ml:cross-sectional-deep-models:ae}{autoencoder} 0.15, the contrastive encoder 0.05.
\item Ten features fitted on two and a half minutes of mostly flat flow are noise; with more data they carry the
microstructure knowledge the networks must learn.
\item Ridge on the raw window 0.45, the predictive probe 0.40, fine-tuned 0.46, from scratch 0.38, the \omterm{def:ml:cross-sectional-deep-models:ae}{autoencoder} and
contrastive probes 0.31, the hand-made features 0.55.
\item \omterm{def:ml:representation-learning:ssl}{Fine-tuning} adapts the encoder's code to the task once there are enough labels to do so safely; a frozen code keeps
only what the pretext kept.
\item $0.46 - 0.38 = 0.08$ of IC.
\item The curves to converge: with many labels, \omterm{def:ml:representation-learning:ssl}{pre-training}'s advantage fades and ridge on raw or hand-made inputs catch
up.
\item \emph{Named result}: below about 1\,000 labelled windows (eight minutes) the predictive probe beats every raw-input
model trained on labels alone: 0.22 against 0.18 from scratch at 300, 0.29 against 0.20 at 1\,000; with 4\,444 only the
fine-tuned encoder keeps level with ridge.
\item When labels are scarce and unlabelled data plentiful, and when no one has written down the features (a new market,
a new data type).
\item From what should not change a forecast: small measurement noise, shifts by a snapshot, rescaling of deep levels; test
them with a probe, not by taste.
\item The test period and anything after it.
\item Pool sessions of many instruments, with per-instrument normalisation, and probe on each.
\item How much of the task is linearly available in the representation, independently of how a downstream model is
trained.
\item It is a very large self-supervised forecaster of text, used through probes, prompts or \omterm{def:ml:representation-learning:ssl}{fine-tuning}.
\item Being close to the task, so that what it keeps is what the task needs.
\end{enumerate}
\end{solution}

\begin{solution}{iq:ml:representation-learning:1}
A \omterm{def:ml:representation-learning:probe}{linear probe} fits a linear model on frozen codes: cheap, measures what the representation holds, safe with few labels.
\omterm{def:ml:representation-learning:ssl}{Fine-tuning} trains the encoder too: can adapt the representation, needs more labels and care (small \omterm{def:ml:trees-and-boosting:boosting}{learning rate}, \omterm{def:ml:trees-and-boosting:early}{early
stopping}). Probe first; fine-tune when there are enough labels.
\iqlookfor{the trade-off between adaptation and \omterm{def:ml:why-financial-machine-learning-is-different:generalisation}{overfitting}.}
\end{solution}

\begin{solution}{iq:ml:representation-learning:2}
Train codes so that two views of one example are close and views of different examples far (InfoNCE). For intraday
prices: small noise, small time shifts, masking of some features, rescaling of volumes; never transformations that change
direction or future outcome (reversing time, flipping sides).
\iqlookfor{the objective and augmentations that respect what matters.}
\end{solution}

\begin{solution}{iq:ml:representation-learning:3}
Pre-train an encoder on the ten years with a self-supervised pretext close to the task (forecasting the next seconds or
days of the stream), then probe and possibly fine-tune on the three months; compare on held-out labelled data with a model
trained from scratch.
\iqlookfor{a pretext aligned with the task and an honest comparison.}
\end{solution}

\begin{solution}{iq:ml:representation-learning:4}
It keeps what varies most in the input (PCA for a linear \omterm{def:ml:cross-sectional-deep-models:ae}{autoencoder}), and the predictive part of the input may be a
low-variance direction the code drops.
\iqlookfor{variance is not relevance.}
\end{solution}

\begin{solution}{iq:ml:representation-learning:5}
\omterm{def:ml:representation-learning:probe}{Linear probes} on held-out data for several tasks, and the learning curve in labelled data against a raw-input baseline.
\iqlookfor{probes and learning curves.}
\end{solution}

\begin{solution}{iq:ml:representation-learning:6}
The reconstruction $BAX$ has rank at most $k$; by Eckart--Young the best rank-$k$ Frobenius approximation of $X$ is the
truncated SVD, i.e.\ the projection onto the top $k$ right singular vectors, the principal directions; choosing $A$ as the
projection's coordinates and $B$ as its basis attains it.
\iqlookfor{the rank argument and Eckart--Young.}
\end{solution}
